% Generated draft: requires mathematical and source review.
\documentclass[11pt]{article}
\usepackage[margin=25mm]{geometry}
\usepackage{fontspec}
\setmainfont{Latin Modern Roman}
\usepackage{amsmath,amssymb,mathtools,unicode-math}
\setmathfont{Latin Modern Math}
\usepackage{enumitem,needspace,xurl,hyperref}
\setlength{\parindent}{0pt}
\setlength{\parskip}{5pt}
\setlength{\emergencystretch}{3em}
\newcommand{\N}{\mathbb{N}}
\newcommand{\Z}{\mathbb{Z}}
\newcommand{\Q}{\mathbb{Q}}
\newcommand{\R}{\mathbb{R}}
\newcommand{\C}{\mathbb{C}}
\newcommand{\F}{\mathbb{F}}
\newcommand{\A}{\mathbb{A}}
\newcommand{\PP}{\mathbb{P}}
\newcommand{\E}{\mathbb{E}}
\newcommand{\eps}{\varepsilon}
\newcommand{\dd}{\,\mathrm{d}}
\begin{document}
\section*{UnsolvedMath ID: 6800005}
\textbf{Source problem code:} catalogue entry
\subsection{Definitions and mathematical statement}
An area-minimizing \(m\)-current \(T\) minimizes mass among the admissible currents with the same boundary, in the global or local class specified by its setting. The supplied excerpt does not specify that setting, coefficients, multiplicities, or the definition of \(\operatorname{Sing}_b(T)\). In particular the subscript \(b\) must not be silently identified with branching or boundary singularities. Let \(\mathcal H^s\) and \(\dim_{\mathrm H}\) denote Hausdorff measure and dimension. De Lellis poses the following questions:
\begin{enumerate}
\item Is \(\mathcal H^{m-1}(\operatorname{Sing}_b(T))=0\)?
\item If the preceding answer is yes, must \(\dim_{\mathrm H}\operatorname{Sing}_b(T)\leq m-2\)?
\item Can the example referred to in the missing citation be modified so that its accumulating singularities \(P_k\in\operatorname{Sing}_b(T)\) are of branching type? The source suggests that a positive answer might enable a Cantor-type construction contradicting the dimension bound in item 2. It does not assert that such a construction already follows or exists.
\item White's conjecture, as reported here, says that a 2-dimensional area-minimizing current spanning a real analytic closed curve has finitely many singularities in total, including both interior and boundary singularities, and hence finite topological type.
\item Does there exist a smooth closed curve \(\Gamma\subset\R^{2+n}\) bounding an area-minimizing 2-current \(T\) with infinite topology? In an oriented multiplicity-one integral-current interpretation, the boundary condition is \(\partial T=\lBrack\Gamma\rBrack\), but those coefficient and multiplicity conventions are not given. The source notes that the existing example in another missing citation is topologically a disk.
\end{enumerate}
The definition of the relevant singular set and the referenced accumulating example are indispensable for items 1--3. Item 3 itself distinguishes belonging to \(\operatorname{Sing}_b(T)\) from being of branching type. The meaning of the topology of a possibly singular current, and the allowed codimension \(n\) in item 5, also require clarification. Infinite topology is not automatically the same condition as infinitely many singular points.
\subsection{Short English statement}
Ask for measure and dimension bounds on the specified singular set of an area minimizer, and for accumulating branching examples. Also retain White's finite-singularity conjecture for analytic boundary curves and the contrasting question of a smooth boundary curve spanning a minimizer with infinite topology.
\subsection{Sources}
\begin{itemize}
\item Local UnsolvedMath registry, permanent ID 6800005, source code catalogue entry; attributed to Camillo De Lellis, Institute for Advanced Study, Princeton; proposed year 2018; update 2026-08-06.
\item Morgan and Pansu, \emph{A List of Open Problems in Differential Geometry} (2018), Question environment 5, ``Branch points of area-minimizing surfaces'': \url{https://www.imo.universite-paris-saclay.fr/~pansu/problems_MTDG.tex}. Registry access date 2026-07-29; source-TeX extraction; public source on a coauthor's publication page. Current status and redistribution terms require review.
\item White is credited with the analytic-boundary conjecture, but the citation following ``White in'' is absent. Citations following ``example as in'' and ``example in'' are also absent.
\item The supplied summary mentions De Lellis--Spadaro--Marchese regularity and singular-set work without a precise reference, and reports the questions as unresolved as of August 2026. These are unverified registry assessments.
\end{itemize}
\end{document}
